\clearpage
\VerdierPhysicalPageAnchor{0077}%
\VerdierSemanticAnchor{ch1-2-4-3-simple-complex-morphisms}
于是，取单复形后，得到两个函子态射：

\VerdierSemanticAnchor{ch1-2-4-3-eq2}
\begin{equation}
  \begin{aligned}
    \operatorname{comp}_{\Pi}h\colon
      \operatorname{comp}_{\Pi}F
      &\longrightarrow \operatorname{comp}_{\Pi}F',\\
    \operatorname{comp}_{\Sigma}h\colon
      \operatorname{comp}_{\Sigma}F
      &\longrightarrow \operatorname{comp}_{\Sigma}F'.
  \end{aligned}
  \tag{2.4.3.2}
\end{equation}

\VerdierSemanticAnchor{ch1-2-4-3-composition}
若 $h$ 与 $h'$ 是多加法函子的两个可复合态射，则有：
\begin{equation*}
  \begin{aligned}
    \operatorname{comp}_{\Pi}hh'
      &=\operatorname{comp}_{\Pi}h\,\operatorname{comp}_{\Pi}h',\\
    \operatorname{comp}_{\Sigma}hh'
      &=\operatorname{comp}_{\Sigma}h\,\operatorname{comp}_{\Sigma}h'.
  \end{aligned}
\end{equation*}

\VerdierSemanticAnchor{ch1-2-4-3-sign-rule-ii}
态射 $\operatorname{comp}_{\Pi}h$（相应地，
$\operatorname{comp}_{\Sigma}h$）是遵循符号规则 II (1.6.6) 的
$\mathbb Z[n]$-函子态射。

\VerdierSemanticAnchor{ch1-2-4-4}
\noindent\textbf{2.4.4.}\quad
当函子 $F$ 在某些变量上反变，即对 $i\in J\subset[n]$ 有等式
$C_i=(C_i')^\circ$ 时，模去 (2.4.1) 中所作的识别，函子
$\operatorname{comp}_{\Pi}F$ 与 $\operatorname{comp}_{\Sigma}F$
在相应变量上反变。

\VerdierSemanticAnchor{ch1-2-5}
\subsection{同伦}

\VerdierSemanticAnchor{ch1-2-5-1}
\noindent\textbf{2.5.1.}\quad
设 $X$ 与 $Y$ 是 $C$ 中的两个 $n$-重复形，
$f,g\colon X\rightrightarrows Y$ 是两个复形态射。连接 $f$ 与 $g$
的一个\emph{同伦}，是如下的一族态射

\VerdierSemanticAnchor{ch1-2-5-1-family}
\begin{equation*}
  a=\bigl(a^{\sigma,i}\colon X^{\sigma+e_i}\longrightarrow Y^\sigma,
    \ \sigma\in\mathbb Z[n],\ i\in[n]\bigr)
\end{equation*}
使得：

\VerdierSemanticAnchor{ch1-2-5-1-eq1}
\begin{equation}
  \begin{aligned}
    a^{\sigma+e_i-e_j,j}d_X^{\sigma,i}
      &=d_Y^{\sigma-e_j,i}a^{\sigma-e_j,j},
      &&\forall\sigma\in\mathbb Z[n],\
        \ \forall i,j\in[n],\ i\ne j,\\
    g^\sigma-f^\sigma
      &=\sum_{i\in[n]}\Bigl(
          a^{\sigma,i}d_X^{\sigma,i}
          +d_Y^{\sigma-e_i,i}a^{\sigma-e_i,i}\Bigr),
      &&\forall\sigma\in\mathbb Z[n].
  \end{aligned}
  \tag{2.5.1.1}
\end{equation}

\VerdierSemanticAnchor{ch1-2-5-1-notation}
用记号 $a\colon f\rightsquigarrow g$ 表示连接 $f$ 与 $g$ 的同伦 $a$。

\VerdierSemanticAnchor{ch1-2-5-2}
\noindent\textbf{2.5.2.}\quad
设 $f,g,h\colon X\longrightarrow Y$ 是三个态射，
$a\colon f\rightsquigarrow g$ 与 $b\colon g\rightsquigarrow h$ 是两个同伦。
用 $b+a\colon f\rightsquigarrow h$ 表示同伦

\VerdierSemanticAnchor{ch1-2-5-2-sum}
\begin{equation*}
  (b+a)^{\sigma,i}=b^{\sigma,i}+a^{\sigma,i},
  \qquad \sigma\in\mathbb Z[n],\ i\in[n].
\end{equation*}

\VerdierSemanticAnchor{ch1-2-5-3}
\noindent\textbf{2.5.3.}\quad
设 $f,g\colon X\rightrightarrows Y$ 是两个复形态射，
$a\colon f\rightsquigarrow g$ 是一个同伦，而
$u\colon Y\longrightarrow Y'$（相应地，$v\colon X'\longrightarrow X$）
是一个复形态射。用
$u\star a\colon uf\rightsquigarrow ug$（相应地，
$a\star v\colon fv\rightsquigarrow gv$）表示同伦

\VerdierSemanticAnchor{ch1-2-5-3-whiskering}
\begin{equation*}
  (u\star a)^{\sigma,i}=u^\sigma a^{\sigma,i}
  \qquad
  \text{（相应地，}(a\star v)^{\sigma,i}
    =a^{\sigma,i}v^{\sigma+e_i}\text{）}.
\end{equation*}
